% ===== §7 Objections and Limits =====
\section{Objections and Limits}
\label{sec:objections}

\noindent
Four objections press hardest, and the section meets each not to dispose of it cheaply but to let it fix the boundary of the claim. The second, the autonomy objection, is the one on which the paper is most exposed; it receives the fullest treatment, drawing on the literature that has thought longest about wrongs done with the victim's consent.

\subsection{That it counsels the able to be less}
\label{sub:obj-less}
The first objection is the one the construction was built to answer, and it returns for final disposal. If effacement is the harm one party's superior generation does another, does the theory not simply counsel the able to do less---to dim their gifts so that others may shine by comparison? It does not, and the answer is now exact rather than merely intuitive. The criterion of \xref{sec:criterion} is not the \emph{rate} of one party's generation but the \emph{capture} of the other's frequency, and capture is a function of coupling, not of speed: the locking threshold $K_{\text{lock}}=|\omega_s-\omega_w|/2$ binds $K$, not $\omega_s$. The able party is asked not to generate less but to cease entraining---to release the traction by which its rhythm overwrites another's, while keeping its own rhythm entire. The theory counsels decoupling, not diminishment, and \xref{sec:reticence} showed these to be not merely different but opposed: deceleration keeps the strong pole the pace-setter and adds the second harm of a sacrifice owed, where decoupling restores the other's frequency at no cost to the stronger's own. The objection mistakes a counsel of non-determination for a counsel of self-reduction; the model is precisely what tells them apart.

\subsection{The autonomy objection}
\label{sub:obj-autonomy}
The second objection is the one on which the paper is most exposed, and it neither evades nor softens it. If the effaced party consents to her effacement---endorses it, is grateful for it, would choose it again---then by what right does the theory pronounce her wronged, and pronounce it, apparently, over her own sincere judgement of her own life? To do so seems the very arrogance the paper condemns elsewhere: the substitution of the theorist's verdict for the subject's self-understanding.

The objection cannot be answered by denying its premise, for the premise is true and the paper has insisted on it: the effaced party often does consent, and her consent is not feigned. It must be answered by showing that consent, here, does not do the exculpatory work the objection asks of it---and the resources for showing this are already developed, in the literature on \emph{adaptive preferences}. Elster named the structure: a ``blind psychic process operating behind the back of the person,'' by which an agent revises her preferences to fit her diminished opportunities without awareness or control \citep{elster1983}. Nussbaum gave it ethical weight through cases like Vasanti's, the woman who came to regard her husband's abuse as part of a woman's lot, her preference formed by lifelong habituation to constrained options \citep{nussbaum2001}. The lesson the literature draws is precisely the one the autonomy objection overlooks: that a preference can be sincerely held, reflectively endorsed, and \emph{still} fail to exculpate the conditions that formed it, because the endorsement itself may be an artefact of those conditions. The effaced party's gratitude is consistent with her effacement in exactly this way---a generativity foreclosed long enough comes to feel like a preference satisfied.

But the adaptive-preference literature also carries a warning the paper must heed, and heeding it is what keeps the reply from becoming the arrogance the objection feared. Narayan and others have cautioned against treating women who endorse self-subordinating arrangements as mere ``dupes'' or ``prisoners of patriarchy,'' urging that their choices be read, often, as competent bargaining within hard constraints rather than as deformed desire \citep{narayan2002}. The caution is decisive here, and it is why the paper's reply has two parts that must be held together. The first: consent does not exculpate, because effacement operates at the level of being and consent at the level of the will, and a will habituated to its own foreclosure can sincerely endorse it. The second, equally firm: that consent does not exculpate gives \emph{no one} licence to override the consenting party, to declare her deceived, or to manage her toward a generativity she has not asked to recover. The wrong is real under consent; \emph{and} the authority over the claim remains hers, by the second-personal ground of \xref{sub:eth-ground} and the reflexive constraint of \xref{sub:eth-paternalism}. The theory thus occupies a position narrower and more disciplined than the objection supposes available: it holds that a wrong can be done with sincere consent, while denying that this holding authorises any intervention over the head of the one who consents. It diagnoses a harm; it does not license a rescue.

\begin{claim}[Consent does not exculpate, and does not authorise]
The autonomy objection is answered by a position with two inseparable halves. \emph{Consent does not exculpate effacement}: as the literature on adaptive preferences establishes, a preference sincerely held and reflectively endorsed may still be an artefact of the very foreclosure it endorses, so that the effaced party's gratitude is consistent with her being wronged. \emph{And the failure of consent to exculpate authorises no intervention}: the second-personal warrant locates the claim's authority in the effaced party herself, so that to override her, or to manage her recovery, would be a fresh effacement. The paper holds both: the wrong survives consent, and no one but the wronged has standing to act on its being a wrong. To collapse either half---to say consent cures the harm, or that the harm licenses rescue---is to leave the narrow ground on which the diagnosis is both true and non-arrogant.
\end{claim}

\subsection{The asymmetry-of-capacity objection}
\label{sub:obj-capacity}
The third objection holds that sometimes one party simply \emph{can} generate more---is more gifted, more practised, better suited to the problems at hand---and asks whether the theory does not then penalise a mere difference of capacity, demanding an artificial equalisation that wrongs the abler party in the name of sparing the other. The reply distinguishes two things the objection runs together: \emph{generating more} and \emph{occupying the domain of the other's generation}. A more capable party may generate a great deal without effacing, provided the relational field is constituted so that the other's own frequency has a domain in which to run; conversely, a less capable party is not effaced by the mere existence of someone abler, but by the foreclosure of the terrain on which her own---perhaps quite different---capability would operate. This is the substance of the third channel of \xref{sub:ret-channels}: the remedy for the harm of completeness is not the abler party's self-diminishment but the constitution of the shared field to include the terrain where the other's frequency is highest. Reticence is not the equalisation of capacities and asks no one to feign incapacity; it asks only that the field not be drawn solely around one party's strengths. The difference of capacity is not the harm; the totality of the occupation is.

\subsection{The limit of the criterion}
\label{sub:obj-overclaim}
The fourth is less an objection than a limit the paper imposes on itself, and states plainly to forestall misuse. The formal criterion of \xref{sec:criterion} detects a \emph{structural risk}; it locates the conditions under which effacement occurs. It is not an instrument for the verdict on an actual relation, and still less for the judgement of persons. Real relations are not two oscillators; the model abstracts severely, and its worth is conceptual clarification, not measurement. To read off from it that a given partner stands convicted of effacing another would mistake an illumination for a tribunal---and would, moreover, violate the very constraint of \xref{sub:eth-paternalism}, since to convict on another's behalf is itself to determine her domain. What the criterion warrants is a cultivated reticence: a sensitised attention to the possibility of the harm, a disposition to look for it where it hides behind devotion. It is a lamp held to a dark region, not a gavel; and a lamp is what a harm this well-concealed requires, for what announces itself by no transgression can be sought only by a perception taught where to look.
